\documentclass[10pt,a4paper]{amsart}
\usepackage{amsmath,amssymb,amsthm,graphicx}
% Retained minimal preamble for the Stacks Project chapter more-morphisms.tex (More on Morphisms).
% Verbatim excerpt of the chapter preamble: Xy-pic setup (native preamble lines 14-18) and the native
% theorem declarations with the mathematical macros (native preamble lines 154-217). Copied mechanically;
% nothing added, removed or redefined. The wrapper class and the amsmath/amssymb/amsthm/graphicx base
% packages are supplied by the assembler, so the class/externaldocument part of the native preamble is omitted.
% For commutative diagrams we use Xy-pic
\usepackage[all]{xy}

% We use 2cell for 2-commutative diagrams.
\xyoption{2cell}
% Theorem environments.
%
\theoremstyle{plain}
\newtheorem{theorem}[subsection]{Theorem}
\newtheorem{proposition}[subsection]{Proposition}
\newtheorem{lemma}[subsection]{Lemma}

\theoremstyle{definition}
\newtheorem{definition}[subsection]{Definition}
\newtheorem{example}[subsection]{Example}
\newtheorem{exercise}[subsection]{Exercise}
\newtheorem{situation}[subsection]{Situation}

\theoremstyle{remark}
\newtheorem{remark}[subsection]{Remark}
\newtheorem{remarks}[subsection]{Remarks}

\numberwithin{equation}{subsection}

% Macros
%
\def\lim{\mathop{\mathrm{lim}}\nolimits}
\def\colim{\mathop{\mathrm{colim}}\nolimits}
\def\Spec{\mathop{\mathrm{Spec}}}
\def\Hom{\mathop{\mathrm{Hom}}\nolimits}
\def\Ext{\mathop{\mathrm{Ext}}\nolimits}
\def\SheafHom{\mathop{\mathcal{H}\!\mathit{om}}\nolimits}
\def\SheafExt{\mathop{\mathcal{E}\!\mathit{xt}}\nolimits}
\def\Sch{\mathit{Sch}}
\def\Mor{\mathop{\mathrm{Mor}}\nolimits}
\def\Ob{\mathop{\mathrm{Ob}}\nolimits}
\def\Sh{\mathop{\mathit{Sh}}\nolimits}
\def\NL{\mathop{N\!L}\nolimits}
\def\CH{\mathop{\mathrm{CH}}\nolimits}
\def\proetale{{pro\text{-}\acute{e}tale}}
\def\etale{{\acute{e}tale}}
\def\QCoh{\mathit{QCoh}}
\def\Ker{\mathop{\mathrm{Ker}}}
\def\Im{\mathop{\mathrm{Im}}}
\def\Coker{\mathop{\mathrm{Coker}}}
\def\Coim{\mathop{\mathrm{Coim}}}

% Boxtimes
%
\DeclareMathSymbol{\boxtimes}{\mathbin}{AMSa}{"02}

%
% Macros for moduli stacks/spaces
%
\def\QCohstack{\mathcal{QC}\!\mathit{oh}}
\def\Cohstack{\mathcal{C}\!\mathit{oh}}
\def\Spacesstack{\mathcal{S}\!\mathit{paces}}
\def\Quotfunctor{\mathrm{Quot}}
\def\Hilbfunctor{\mathrm{Hilb}}
\def\Curvesstack{\mathcal{C}\!\mathit{urves}}
\def\Polarizedstack{\mathcal{P}\!\mathit{olarized}}
\def\Complexesstack{\mathcal{C}\!\mathit{omplexes}}
% \Pic is the operator that assigns to X its picard group, usage \Pic(X)
% \Picardstack_{X/B} denotes the Picard stack of X over B
% \Picardfunctor_{X/B} denotes the Picard functor of X over B
\def\Pic{\mathop{\mathrm{Pic}}\nolimits}
\def\Picardstack{\mathcal{P}\!\mathit{ic}}
\def\Picardfunctor{\mathrm{Pic}}
\def\Deformationcategory{\mathcal{D}\!\mathit{ef}}
\usepackage{cleveref}
\crefname{originallemma}{original lemma}{original lemmas}
\Crefname{originallemma}{Original lemma}{Original lemmas}
\makeatletter
\newlabel{lemma-etale-splits-off-quasi-finite-part-technical-variant}{{02LO}{source}{}{sourcepointer-1}{}}
\newlabel{lemma-etale-splits-off-quasi-finite-part-technical-variant@cref}{{[originallemma][0][]02LO}{[1][0][]source}{}{}{}}
\makeatother
\crefname{originallemma}{original lemma}{original lemmas}
\Crefname{originallemma}{Original lemma}{Original lemmas}
\makeatletter
\newlabel{lemma-weighting-check-after-etale-base-change}{{0F39}{source}{}{sourcepointer-2}{}}
\newlabel{lemma-weighting-check-after-etale-base-change@cref}{{[originallemma][0][]0F39}{[1][0][]source}{}{}{}}
\makeatother
\title{The integral of a positive weighting is lower semi-continuous}
\author{The Stacks Project Authors}
\date{stacks-project:more-morphisms.tex:local-snapshot-20261011; GNU FDL 1.2 or later}
\begin{document}
\maketitle
\noindent\textit{Editorial note.} Counted span is the author's own proof environment, source lines 22702-22728, immediately following the statement at 22694-22700.\par\medskip
\noindent\textit{Editorial note.} The proof cites the etale splitting lemma and uses the defining local constancy of a weighting; both are navigation pointers into the same chapter and are not re-counted.\par\medskip
\noindent\textit{Editorial note.} The displayed chain of equalities for the integral over U is retained verbatim, including the textstyle integrals.\par\medskip
\noindent\textit{Allowed standard prerequisites (editorial).}\begin{itemize}
\item Etale splitting off of the quasi-finite part. chapter lemma lemma-etale-splits-off-quasi-finite-part-technical-variant, cited at proof lines 22704-22706. Gives the etale neighbourhood and the decomposition of U x\_Y X into W and the V\_\{i,j\} used to write the integral as a finite sum of locally constant pieces.
\item Criterion for a weighting after etale base change. chapter lemma lemma-weighting-check-after-etale-base-change, cited via the definition of the integral. Explains why each integral over V\_\{i,j\} of the restricted weighting is locally constant, hence constant after shrinking U.
\item Base change compatibility of the integral. used at proof lines 22725-22726. Transfers the lower bound obtained over the etale neighbourhood back to Y.
\end{itemize}\medskip
\ifdefined\hypertarget\hypertarget{sourcepointer-1}{}\fi
\noindent\textit{Source reference.} Original lemma 02LO: sources/stacks-more-morphisms/pinned-original-source.tex, source lines 12136--12136. Navigation pointer only: the counted proof cites this statement by \textbackslash{}ref as a standing input; the deduction the counted proof has to make is carried out inside the counted span, and no part of that deduction is delegated to this pointer.\par\medskip
\ifdefined\hypertarget\hypertarget{sourcepointer-2}{}\fi
\noindent\textit{Source reference.} Original lemma 0F39: sources/stacks-more-morphisms/pinned-original-source.tex, source lines 22111--22111. Navigation pointer only: the counted proof cites this statement by \textbackslash{}ref as a standing input; the deduction the counted proof has to make is carried out inside the counted span, and no part of that deduction is delegated to this pointer.\par\medskip
\begin{lemma}
\label{lemma-semicontinuous-int-w}
Let $f : X \to Y$ be a separated, locally quasi-finite morphism
with finite fibres.
Let $w : X \to \mathbf{Z}_{> 0}$ be a positive weighting of $f$.
Then $\int_f w$ is lower semi-continuous.
\end{lemma}

\begin{proof}
Let $y \in Y$. Let $x_1, \ldots, x_r \in X$ be the points lying over $y$.
Apply
Lemma \ref{lemma-etale-splits-off-quasi-finite-part-technical-variant}
to get an \'etale neighbourhood $(U, u) \to (Y, y)$ and a decomposition
$$
U \times_Y X =
W \amalg
\ \coprod\nolimits_{i = 1, \ldots, n}
\ \coprod\nolimits_{j = 1, \ldots, m_i}
V_{i, j}
$$
as in locus citatus. Observe that $(\int_f w)(y) = \sum w(v_{i, j})$
where $w(v_{i, j}) = w(x_i)$. Since $\int_{V_{i, j} \to U} w|_{V_{i, j}}$
is locally constant by definition, we may after shrinking $U$
assume these functions are constant with value $w(v_{i, j})$.
We conclude that
$$
\textstyle{\int}_{U \times_Y X \to U} w|_{U \times_Y X} =
\textstyle{\int}_{W \to U} w|_W +
\sum \textstyle{\int}_{V_{i, j} \to U} w|_{V_{i, j}} =
\textstyle{\int}_{W \to U} w|_W + (\int_f w)(y)
$$
This is $\geq (\int_f w)(y)$ and we conclude because $U \to Y$
is open and formation of the integral commutes with base change
(Lemma \ref{lemma-weighting-check-after-etale-base-change}).
\end{proof}

\end{document}
